\documentclass[11pt,a4paper,oneside]{book}
\usepackage{mathscommun}
\hypersetup{pdftitle={Mathématiques de troisième - Guide, couverture et révisions}}
\begin{document}
\titrecollection{Guide, couverture et révisions}{Le mode d'emploi de la collection.}{Périmètre officiel, correspondance des notions, progression, statut des démonstrations, automatismes et organisation des sources LaTeX.}
\frontmatter\tableofcontents\mainmatter
\chapter{Périmètre et sources}
\section{Le niveau choisi}
Ce dossier vise la \textbf{troisième du programme français applicable en 2026--2027}. Il s'appuie sur le programme du cycle 4 publié en 2020 et les attendus de fin d'année de troisième. Il reprend aussi les acquis du collège nécessaires pour rendre le cours autonome.

Un nouveau programme de mathématiques du cycle 4 a été publié en mars 2026. Son application est progressive : cinquième en septembre 2026, quatrième en septembre 2027, troisième en septembre 2028. Il ne remplace donc pas encore le programme de troisième visé ici. Une édition destinée à 2028 devra faire l'objet d'une vérification spécifique ; ce dossier ne revendique pas cette couverture future.

Les preuves et exercices sont rédigés pour ce dossier. Les textes officiels servent à fixer le périmètre, pas à fournir des exercices recopiés.
\section{Références officielles consultées le 15 octobre 2026}
\begin{enumerate}[label={S\arabic*.}]
\item Ministère de l'Éducation nationale, \emph{Programme d'enseignement du cycle des approfondissements (cycle 4)}, annexe 3 de l'arrêté du 17 juillet 2020, BO n°31 du 30 juillet 2020. Consulter la partie Mathématiques, thèmes A à E.\\
\href{https://www.education.gouv.fr/bo/20/Hebdo31/MENE2018714A.htm}{Page du Bulletin officiel}\\
\href{https://cache.media.education.gouv.fr/file/31/89/1/ensel714_annexe3_1312891.pdf}{Annexe 3 : programme du cycle 4 (PDF)}
\item Éduscol, \emph{Mathématiques : attendus de fin d'année de troisième}, document officiel de 11 pages.\\
\href{https://eduscol.education.gouv.fr/sites/default/files/document/18-maths-3e-attendus-eduscol1114748pdf-74688.pdf}{Attendus de fin d'année de troisième (PDF)}
\item Éduscol, \emph{Ressources d'accompagnement du programme de mathématiques au cycle 4}, mise à jour d'avril 2026, calendrier d'application et accès aux programmes, repères et attendus.\\
\href{https://eduscol.education.gouv.fr/5736/ressources-d-accompagnement-du-programme-de-mathematiques-au-cycle-4}{Ressources et calendrier officiel}
\end{enumerate}
\section{Quatre documents, une même collection}
\begin{center}
\begin{tabularx}{\textwidth}{lX}
\toprule Document & Usage\\\midrule
\texttt{cours\_3e.pdf} & Comprendre les notions dans vingt chapitres, avec preuves et exemples.\\
\texttt{exercices\_3e.pdf} & Chercher : quatre exercices par chapitre et quatre synthèses.\\
\texttt{corriges\_3e.pdf} & Vérifier la démarche, identifier l'erreur et recommencer.\\
\texttt{guide\_3e.pdf} & Contrôler la couverture, organiser le travail, réviser et compiler les sources.\\\bottomrule
\end{tabularx}
\end{center}
\chapter{Table de couverture du programme}
La correspondance ci-dessous croise S1 et S2. Elle regroupe les notions pour rester lisible ; les chapitres comportent leurs définitions, méthodes, exemples et justifications. Les quatre exercices de chaque chapitre constituent l'entraînement associé.

\small
\begin{longtable}{p{.64\textwidth}p{.26\textwidth}}
\toprule Notions et savoir-faire & Chapitres du cours\\\midrule
\endfirsthead
\toprule Notions et savoir-faire & Chapitres du cours\\\midrule\endhead
\multicolumn{2}{l}{\textbf{Nombres et calculs}}\\
Relatifs, opposés, inverses ; différentes écritures & 2\\
Fractions, égalité, ordre, opérations justifiées & 2\\
Repérage, encadrements, calcul exact et approché & 2, 4, 18\\
Puissances entières, notation scientifique, préfixes & 4, 17\\
Carrés parfaits, racines, problèmes numériques & 4, 6, 13\\
Division euclidienne, multiples et diviseurs & 3\\
Critères de divisibilité, nombres premiers jusqu'à 100 & 3\\
Décomposition, fractions irréductibles, périodicité & 3, 19\\
Statuts des lettres, identité, opposé d'expression & 5\\
Distributivités, réduction, factorisation & 5\\
Différence de carrés ; produit nul & 5, 6\\
Premier degré, équations produits, carrés & 6\\
Modélisation et preuves par calcul littéral & 1, 5, 6, 20\\
\midrule
\multicolumn{2}{l}{\textbf{Données, probabilités et fonctions}}\\
Recueil, tableaux, effectifs, fréquences & 10, 20\\
Bâtons, secteurs, histogrammes d'amplitude constante & 10\\
Moyenne, médiane, étendue ; interprétation & 10\\
Issues, événements, probabilités, contraire & 11\\
Une ou deux épreuves ; dénombrement & 11\\
Fréquences, probabilités, simulation & 11, 19, 20\\
Proportionnalité, quatrième proportionnelle, ratios & 7\\
Pourcentages, évolution, coefficient multiplicateur & 7, 9\\
Échelles et proportionnalité en géométrie & 7, 14, 16\\
Fonction : image, antécédent, notations & 8\\
Formule, tableau, courbe, programme de calcul & 8, 20\\
Fonctions linéaires et affines, graphiques, paramètres & 9\\
Phénomènes continus, domaines et modèles & 8, 9, 20\\
\midrule
\multicolumn{2}{l}{\textbf{Grandeurs et mesures}}\\
Périmètres, aires, unités, ordres de grandeur & 12, 17\\
Prisme, cylindre, pyramide, cône, boule & 17\\
Assemblages, retraits, volumes et contenances & 17, 20\\
Grandeurs produit et quotient, conversions & 17\\
Effets des transformations : longueurs, aires, volumes & 16\\
\midrule
\multicolumn{2}{l}{\textbf{Espace et géométrie}}\\
Repères : droite, plan, pavé, coordonnées & 2, 8, 18\\
Sphère : latitude, longitude, grand cercle & 18\\
Solides, perspectives, vues, patrons, sections & 18\\
Angles et parallélisme ; triangles et constructions & 12\\
Hauteurs, médiatrices, somme des angles & 12\\
Inégalité triangulaire, triangles égaux & 12\\
Parallélogrammes et propriétés caractéristiques & 12\\
Pythagore : théorème, réciproque, applications & 13\\
Thalès : emboîtement, papillon, réciproque & 14\\
Triangles semblables et côtés proportionnels & 14\\
Cosinus, sinus, tangente : côtés et angles & 15\\
Symétries, translation, rotation, homothétie & 16\\
Frises, pavages, rosaces et transformations & 16, 19\\
\midrule
\multicolumn{2}{l}{\textbf{Algorithmique et outils}}\\
Algorithmes débranchés, lecture et ordre des instructions & 19\\
Variables, affectations, conditions, événements & 19\\
Boucles bornées et conditionnelles, tests, corrections & 19\\
Blocs personnalisés, motifs, boucles imbriquées & 19\\
Simulations et programmes de calcul & 19\\
Scripts parallèles, interactions, jeu à deux lutins & 19\\
Tableur-grapheur, calculatrice, géométrie dynamique & 20\\
\bottomrule
\end{longtable}
\normalsize
\section{Les six compétences travaillées}
\begin{center}
\begin{tabularx}{\textwidth}{lX}
\toprule Compétence & Mise en pratique dans le dossier\\\midrule
Chercher & Essais, représentations, décomposition d'un problème.\\
Modéliser & Équations, fonctions, grandeurs et situations aléatoires.\\
Représenter & Graphiques, tableaux, figures, patrons, arbres.\\
Raisonner & Preuves, réciproques, contre-exemples, contrôle des hypothèses.\\
Calculer & Exact, mental, approché, littéral et instrumenté.\\
Communiquer & Explications, unités, rédaction structurée, protocoles.\\\bottomrule
\end{tabularx}
\end{center}
\section{Acquis repris et approfondissements}
Les fractions, le calcul avec les relatifs, certaines propriétés géométriques, Pythagore, les statistiques de base et le repérage sont en partie des acquis antérieurs du cycle. Ils sont intégralement repris dans la mesure nécessaire à la troisième.

Les compléments sont signalés dans le cours : notations des ensembles de nombres, preuve d'irrationalité de $\sqrt2$, formalisme PGCD/PPCM et algorithme d'Euclide, carrés de sommes et différences, inéquations, valeurs trigonométriques exactes, aire de sphère et optimisation d'une aire. Ils enrichissent le cours sans devenir des attendus obligatoires du périmètre choisi.
\chapter{Démonstrations : ce qui est prouvé et ce qui est admis}
\section{Le contrat de lecture}
Une définition n'a pas à être prouvée : elle fixe un sens. Une règle de calcul ou une propriété géométrique doit avoir une justification. Le dossier donne des preuves élémentaires lorsqu'elles sont accessibles et explicite les dépendances lorsque l'on s'appuie sur un résultat admis.

Une preuve par un dessin consiste à expliquer un découpage, une superposition ou des égalités valables en général. Mesurer une figure particulière ou exécuter une simulation ne constitue pas une preuve générale.
\section{Principaux résultats démontrés}
\begin{center}
\begin{tabularx}{\textwidth}{lX}
\toprule Chapitres & Preuves et justifications\\\midrule
1--2 & Parité, exemples et contre-exemples ; signes ; égalité et opérations de fractions.\\
3 & Division euclidienne, divisibilité, test jusqu'à la racine, existence de la décomposition.\\
4--6 & Règles de puissances, propriétés de racines, irrationalité de $\sqrt2$, identités, produit nul et équations.\\
7--9 & Ratios, coefficients d'évolution, compositions, variation affine et représentation par une droite.\\
10--11 & Moyenne pondérée, rangs médians, fréquences, dénombrement, complémentarité.\\
12 & Angles du triangle, triangles égaux, isocèle, médiatrice, cercle circonscrit, parallélogrammes.\\
13--15 & Pythagore et réciproque ; Thalès et réciproque ; similitude ; cohérence des rapports trigonométriques.\\
16 & Conservations des transformations ; facteurs de longueurs, aires et volumes.\\
17--18 & Aires polygonales ; volumes sous principes admis ; boule par sections ; patrons et sections.\\
19--20 & Correction de boucles simples, angles de rotation, algorithme de factorisation, aire maximale.\\\bottomrule
\end{tabularx}
\end{center}
\section{Socle admis, sans fausse promesse de preuve absolue}
\begin{enumerate}
\item \textbf{Nombres :} les propriétés de base des opérations, notamment la distributivité sur les réels ; leur ordre ; l'existence d'une racine positive pour tout réel positif. Les conséquences utiles sont démontrées.
\item \textbf{Arithmétique :} l'unicité de la décomposition en facteurs premiers est énoncée et admise ; l'existence est expliquée. Le lemme d'Euclide n'est pas développé.
\item \textbf{Géométrie euclidienne :} existence et unicité des constructions élémentaires, propriétés des déplacements rigides, parallèle unique, caractérisation angulaire des parallèles, segment comme plus court chemin. Les critères de triangles égaux sont justifiés par construction et superposition.
\item \textbf{Mesures :} additivité des aires et volumes, prolongement des mesures rationnelles aux réels et passages à la limite nécessaires aux figures courbes.
\item \textbf{Volumes :} Cavalieri et l'effet d'un étirement dans une seule direction sont admis. À partir d'eux, les formules de pyramide, de cône et de boule sont démontrées.
\item \textbf{Sphère :} sa formule d'aire est un complément admis avec intuition ; la preuve complète n'est pas revendiquée.
\item \textbf{Probabilités :} le modèle probabiliste fournit les probabilités des issues ; le rapprochement des fréquences sur de nombreuses répétitions indépendantes est admis, pas prouvé par simulation.
\end{enumerate}
\section{Éviter les raisonnements circulaires}
La preuve de Pythagore utilise les aires et la somme des angles, pas la trigonométrie. La preuve de Thalès utilise les aires, pas d'abord la similitude. La similitude est ensuite reliée à Thalès, puis elle justifie que les rapports trigonométriques ne dépendent que de l'angle. La formule de la boule est obtenue à partir des sections et du cône.

Les chapitres sont regroupés par thèmes pour faciliter la consultation ; une preuve peut renvoyer à un chapitre ultérieur. Par exemple, la droite représentant une fonction affine s'appuie sur Thalès. Le renvoi est alors explicite.
\chapter{Progression et méthode de travail}
\section{Une progression possible sur une année}
Cette proposition répartit le travail en blocs et laisse une marge pour les reprises. Ce n'est pas un calendrier réglementaire : ajuster le rythme aux acquis réels.
\begin{center}
\small
\begin{tabularx}{\textwidth}{lXX}
\toprule Période indicative & Travail principal & Reprises régulières\\\midrule
Semaines 1--4 & 1, 2, 3 : raisonner, nombres, arithmétique & Automatismes et petites preuves\\
Semaines 5--8 & 4, 5, 12 : puissances, littéral, bases géométriques & Fractions, constructions\\
Semaines 9--12 & 6, 7, 13 : équations, proportions, Pythagore & Signes, unités, calcul mental\\
Semaines 13--16 & 8, 9, 14 : fonctions, affines, Thalès & Équations et graphiques\\
Semaines 17--20 & 10, 11, 15 : données, hasard, trigonométrie & Proportions et triangles\\
Semaines 21--24 & 16, 17, 18 : transformations, mesures, espace & Aires, volumes, conversions\\
Toute l'année & 19, 20 : programmation et outils en lien avec les chapitres & Tests, tableur, géométrie dynamique\\
Fin d'année & Quatre synthèses et remédiation ciblée & Retour aux erreurs persistantes\\\bottomrule
\end{tabularx}
\end{center}
\section{Une séance de 45 à 60 minutes}
\begin{enumerate}
\item Cinq minutes de rappel sans le cours : une définition, une formule avec ses conditions, deux calculs.
\item Quinze minutes pour une notion et sa preuve : identifier chaque étape et ce qu'elle utilise.
\item Vingt à trente minutes d'exercices : commencer seul, chercher une représentation, rédiger.
\item Cinq à dix minutes de correction active : expliquer la cause de l'erreur et refaire une question.
\end{enumerate}
Quelques jours plus tard, refaire un exercice de la même famille sans consulter l'exemple. Une réponse juste obtenue uniquement en imitant la ligne précédente ne garantit pas encore la maîtrise.
\section{Une grille d'autoévaluation}
Pour chaque chapitre, noter : « je reconnais », « je sais faire avec aide » ou « je sais faire seul et expliquer ». Pour valider ce dernier niveau, donner une définition avec un exemple, citer les hypothèses d'une propriété, expliquer une preuve, résoudre un exercice nouveau et corriger une erreur typique.
\chapter{Quarante automatismes avec réponses}
Répondre sans calculatrice, sauf si la question demande une approximation. Ces questions courtes complètent les quatre-vingts exercices ; elles ne remplacent pas les problèmes.
\begin{enumerate}
\item Calculer $-7+12$.
\item Calculer $-3-(-8)$.
\item Donner l'opposé de $-5/3$.
\item Donner l'inverse de $-5/3$.
\item Calculer $2/3+1/6$.
\item Calculer $3/4\div9/8$.
\item Comparer $-3/4$ et $-2/3$.
\item Décomposer $84$ en facteurs premiers.
\item Rendre $45/60$ irréductible.
\item $91$ est-il premier ?
\item Calculer $10^{-4}$.
\item Écrire $0{,}00062$ en notation scientifique.
\item Calculer $2^3\times2^4$.
\item Calculer $\sqrt{121}$.
\item Calculer $\sqrt{(-6)^2}$.
\item Développer $3(2x-5)$.
\item Réduire $4x-3+2x+7$.
\item Factoriser $x^2-25$.
\item Résoudre $5x=12$.
\item Résoudre $x^2=16$.
\item Calculer $15\,\%$ de $80$.
\item Donner le coefficient d'une baisse de $12\,\%$.
\item Un prix soldé de $20\,\%$ vaut $40$ euros : quel était son prix ?
\item Partager $50$ selon le ratio $2:3$.
\item Pour $f(x)=3x-2$, calculer $f(-1)$.
\item Pour la même fonction, chercher l'antécédent de $10$.
\item Donner l'ordonnée à l'origine de $g(x)=-2x+5$.
\item Donner la moyenne de $2,4,9$.
\item Donner la médiane de $1,3,8,10$.
\item Donner l'étendue de $5,6,6,12$.
\item Sur un dé équilibré, quelle est la probabilité d'obtenir plus de $4$ ?
\item Si $P(A)=0{,}35$, calculer $P(\overline A)$.
\item Un triangle a deux angles de $40\degre$ et $65\degre$ : donner le troisième.
\item Calculer l'hypoténuse d'un triangle rectangle de côtés $5$ et $12$.
\item Un agrandissement de facteur $3$ multiplie les aires par combien ?
\item Une réduction de facteur $1/2$ multiplie les volumes par combien ?
\item Convertir $2{,}4\unite L$ en $\unite{cm^3}$.
\item Convertir $54\unite{km/h}$ en $\unite{m/s}$.
\item Pour tracer un triangle équilatéral, de combien tourne-t-on après chaque côté ?
\item Une boucle ajoute $3$ à $s$ cinq fois, depuis $s=2$ : valeur finale ?
\end{enumerate}
\section{Réponses}
\begin{longtable}{r p{.80\textwidth}}
\toprule Question & Réponse et contrôle rapide\\\midrule
\endfirsthead\toprule Question & Réponse\\\midrule\endhead
1 & $5$.\\
2 & $5$.\\
3 & $5/3$ : la somme vaut zéro.\\
4 & $-3/5$ : le produit vaut un.\\
5 & $5/6$.\\
6 & $(3/4)(8/9)=2/3$.\\
7 & $-3/4<-2/3$, car $-9/12<-8/12$.\\
8 & $2^2\times3\times7$.\\
9 & $3/4$.\\
10 & Non : $91=7\times13$.\\
11 & $0{,}0001$.\\
12 & $6{,}2\times10^{-4}$.\\
13 & $2^7=128$.\\
14 & $11$.\\
15 & $6$, pas $-6$.\\
16 & $6x-15$.\\
17 & $6x+4$.\\
18 & $(x-5)(x+5)$.\\
19 & $x=12/5=2{,}4$.\\
20 & $x=-4$ ou $x=4$.\\
21 & $12$.\\
22 & $0{,}88$.\\
23 & $40/0{,}8=50$ euros.\\
24 & $20$ et $30$.\\
25 & $-5$.\\
26 & $x=4$.\\
27 & $5$.\\
28 & $(2+4+9)/3=5$.\\
29 & $(3+8)/2=5{,}5$ selon la convention usuelle.\\
30 & $12-5=7$.\\
31 & $2/6=1/3$.\\
32 & $0{,}65$.\\
33 & $75\degre$.\\
34 & $\sqrt{25+144}=13$.\\
35 & $9$.\\
36 & $1/8$.\\
37 & $2400\unite{cm^3}$.\\
38 & $15\unite{m/s}$.\\
39 & $120\degre$, angle extérieur.\\
40 & $2+5\times3=17$.\\
\bottomrule
\end{longtable}
\chapter{Organisation et compilation des sources}
\section{Le style commun}
Le fichier \texttt{mathscommun.sty}, placé dans le dossier \texttt{commun/}, centralise les polices, marges, couleurs, titres, encadrés, figures et notations. Le cours, les exercices, les corrigés et le guide l'utilisent tous. La mise en page est volontairement sobre, adaptée à l'impression.

La commande \verb|\setniveau{Seconde}| permet de réutiliser le style pour un autre niveau. Le niveau affiché dans les titres et pieds de page change sans modifier le style lui-même. La référence au programme et la date se règlent avec \verb|\renewcommand{\ReferenceProgramme}{...}| et \verb|\renewcommand{\EditionDocument}{...}|.
\section{Les fichiers}
Les sources de ce niveau se trouvent dans \texttt{sources/02-college/04-troisieme/} :
\begin{itemize}
\item \texttt{cours\_3e.tex} assemble les vingt fichiers de \texttt{chapitres/}.
\item \texttt{exercices\_3e.tex} assemble les fichiers de \texttt{exercices/}.
\item \texttt{corriges\_3e.tex} assemble les fichiers de \texttt{corriges/}.
\item \texttt{guide\_3e.tex} est le présent guide.
\end{itemize}
Le dossier \texttt{modeles/} contient le modèle d'un nouveau cours. Le script \texttt{compiler.py}, à la racine du dépôt, compile tous les documents. Les PDF compilés se trouvent dans \texttt{build/pdf/troisieme/} et sont publiés dans les releases GitHub ; les fichiers temporaires restent dans \texttt{build/}.

Les figures sont dessinées en TikZ/PGFPlots : elles restent nettes à l'agrandissement et ne dépendent d'aucun fichier image extérieur. Les textes sont encodés en UTF-8.
\clearpage
\section{Compilation conseillée}
Utiliser \textbf{pdfLaTeX} avec une distribution TeX Live ou MiKTeX récente. Les dépendances principales sont \texttt{babel}, \texttt{amsmath}, \texttt{tcolorbox}, \texttt{tikz}, \texttt{pgfplots}, \texttt{subfiles}, \texttt{listings} et \texttt{hyperref}. Une installation LaTeX complète fournit ces extensions.
Depuis la racine du dépôt, lancer :
\begin{lstlisting}
python3 compiler.py
\end{lstlisting}
Le script rend le style commun accessible à LaTeX, lance trois passages pour stabiliser sommaires et renvois, puis vérifie l'absence de débordements et de références non résolues. Une fois toutes les compilations réussies, il actualise les quatre PDF. Pour un chapitre seul :
\begin{lstlisting}
python3 compiler.py --niveau troisieme --chapitre 14
\end{lstlisting}
Ce PDF ponctuel reste dans \texttt{build/}. Les chapitres utilisent \texttt{subfiles} : conserver le fichier maître et l'arborescence. Pour lancer directement pdfLaTeX, ajouter la racine du dépôt à \texttt{TEXINPUTS} et travailler dans le dossier du document. Le script évite cette configuration manuelle.
\section{Ajouter les niveaux suivants}
Conserver l'unique \texttt{mathscommun.sty} comme référence commune. Créer un dossier par niveau à partir du modèle, changer le niveau et le programme applicable, puis ajouter les chapitres et renseigner \texttt{niveaux.json} à la racine. Pour la première et la terminale, préciser le parcours visé avant de revendiquer une couverture exhaustive : voie générale ou technologique, spécialité et options éventuelles. Tous les ajouts et corrections du dépôt sont proposés par pull request, avec leurs sources et les contrôles de leurs PDF, sans versionner les PDF.
\end{document}
